\documentclass[11pt]{article}

\usepackage[margin=1.08in]{geometry}
\usepackage{amsmath,amssymb,amsthm}
\usepackage{enumitem}
\usepackage[hidelinks]{hyperref}

\newtheorem{theorem}{Theorem}[section]
\newtheorem{proposition}[theorem]{Proposition}
\newtheorem{lemma}[theorem]{Lemma}
\newtheorem{corollary}[theorem]{Corollary}
\theoremstyle{definition}
\newtheorem{definition}[theorem]{Definition}
\theoremstyle{remark}
\newtheorem{remark}[theorem]{Remark}

\newcommand{\C}{\mathbb C}
\newcommand{\A}{\mathbb A}
\newcommand{\Jac}{\operatorname{Jac}}
\newcommand{\wt}{\operatorname{wt}}
\newcommand{\ord}{\operatorname{ord}}

\title{A Connected Cubic-Cover Exclusion\\
at Normal Degrees \((9,6)\)}
\author{Atharva Vaidya}
\date{26 July 2026}

\begin{document}
\maketitle

\begin{abstract}
Let \(F,G\in\C[x,y]\) satisfy
\(\Jac(F,G)=\lambda\in\C^\times\), with
\(\deg_yF=9\) and \(\deg_yG=6\).  The leading Jacobian equation
expresses their leading \(y\)-coefficients, after constant rescaling, as
\(h^3\) and \(h^2\) for one \(h\in\C[x]\).  We prove that no such Keller
pair exists when \(h\) is not a cube in \(\C(x)\).

The proof passes to the connected cyclic cover \(H^3=h\), where the
\(\mu _3\)-character and a triangular approximate-root argument give
\[
 f=(g^{3/2})_+ +j(g^{1/2})_+ .
\]
The remaining bracket equations integrate to four weighted first
integrals.  A centered change of coefficient variables decomposes their
level set into two exact branches.  One branch has zero Jacobian.  On
the other, a two-value pole argument excludes every generic level.  At
the exceptional cuspidal level, a finite-map integrality argument
restores polynomial descent and a final valuation identity gives the
contradiction.

The degree pair \((9,6)\) is also the first pair left by a particular
repackaging of earlier partial- and total-degree criteria; that
arithmetic orientation is prior-theorem bookkeeping, not a new
classification result.  The theorem does not cover the chart in which
\(h\) is a cube, and it does not prove the two-dimensional Jacobian
conjecture.
\end{abstract}

\section{Introduction and scope}

The two-dimensional Jacobian conjecture asks whether a polynomial map
\[
 (F,G):\A^2_\C\longrightarrow\A^2_\C
\]
with nonzero constant Jacobian is necessarily a polynomial
automorphism.  The problem remains open.  Earlier work gives strong
restrictions on the total degrees and Newton data of a possible
counterexample; see Guccione--Guccione--Valqui
\cite{GGVShape}, Guccione--Guccione--Horruitiner--Valqui
\cite{GGHVApprox,GGHVAlgorithms}, and the references therein.

This paper studies a different, sharply delimited coefficient chart.
Write
\[
 \deg_yF=9,\qquad \deg_yG=6.
\]
The leading Jacobian equation forces the leading coefficients to be
coprime powers of one polynomial:
\[
 [y^9]F=\alpha h^3,\qquad [y^6]G=\beta h^2,
 \qquad \alpha,\beta\in\C^\times.
\]
The main theorem excludes the case in which \(h\) is not a cube in the
base rational-function field.  Equivalently, it excludes the chart on
which adjoining \(H\) with \(H^3=h\) gives a connected cubic cover.

The connectedness hypothesis is essential to the argument.  It kills
the six nontrivial-character constants that otherwise accompany the
invariant constant \(j\) in the approximate-root expansion.  If \(h\)
is a cube, including when \(h\) is a nonzero constant, those six extra
constants survive and the lower coefficient system is different.
Section~\ref{sec:boundary} records this exact logical boundary.

\paragraph{Relation to earlier results.}
Moskowicz's partial-degree criteria \cite{Moskowicz} and the classical
total-degree consequences used there give the arithmetic orientation
in Section~\ref{sec:arithmetic}.  The argument is included to explain
why \((9,6)\), rather than \((8,6)\), is the first pair not removed by
that particular package of criteria.  It is a repackaging of prior
theorems plus elementary integer enumeration and is not claimed as a
new degree theorem.  The proof avoids the simultaneous-prime assertion
in Lemma 2.3 of \cite{Moskowicz}: in the present common-leading-factor
chart the two reduced progressions coincide, so ordinary Dirichlet
suffices.

Approximate roots and intersection constraints for Jacobian pairs have
been developed extensively, including in \cite{GGHVApprox}; broader
minimal-counterexample and algorithmic frameworks appear in
\cite{GGVShape,GGHVAlgorithms}.  Kara\'s's weighted-bidegree
classification \cite{Karas} provides useful automorphism-side
orientation, but it is not used in the connected-chart exclusion.
Within a scoped audit of these five sources, no prior theorem was
located that states the connected noncube-\(h\), normal-degree
\((9,6)\) exclusion proved here.  This is an audit report, not an
absolute claim of priority: the literature search was not exhaustive,
and independent verification remains necessary before dissemination.

\paragraph{What is not claimed.}
The theorem neither constructs a counterexample nor proves the plane
Jacobian conjecture.  It does not eliminate all Keller pairs of normal
degrees \((9,6)\), because the cube-\(h\) chart remains.  It also does
not assert that every hypothetical counterexample can be transformed
into the present chart.

\section{Statement and arithmetic orientation}

\subsection{The connected-chart theorem}

We use
\[
 [U,V]_{x,y}=U_xV_y-U_yV_x .
\]

\begin{theorem}[Connected cubic-cover exclusion]
\label{thm:main}
Let \(F,G\in\C[x,y]\) satisfy
\[
 [F,G]_{x,y}=\lambda\in\C^\times,\qquad
 \deg_yF=9,\qquad \deg_yG=6.
\]
Suppose their leading \(y\)-coefficients are
\[
 [y^9]F=\alpha h^3,\qquad [y^6]G=\beta h^2
\]
for \(h\in\C[x]\setminus\{0\}\) and
\(\alpha,\beta\in\C^\times\).  Then \(h\) is a cube in \(\C(x)\).
\end{theorem}

Since \(\C\) is algebraically closed, constant rescalings of \(F,G\)
absorb \(\alpha,\beta\).  We therefore prove the theorem with leading
coefficients \(h^3,h^2\).

\subsection{A prior-theorem arithmetic orientation}
\label{sec:arithmetic}

The following proposition is recorded only to locate the chart.  Its
automorphism criteria are consequences of the results collected by
Moskowicz \cite[Theorems 1.1, 1.2, 2.6, and 2.7]{Moskowicz}.

\begin{proposition}[First survivor of the cited arithmetic sieve]
\label{prop:orientation}
Let \(m=\deg_yF\geq n=\deg_yG>0\), and apply:
\begin{enumerate}[label=\textup{(\roman*)}]
 \item the leading-coefficient relation;
 \item constant target reductions when \(m=n\), and the leading-power
 target shear when one reduced exponent is one;
 \item the cited one-coordinate partial-degree criterion; and
 \item the cited total-degree criteria after a sufficiently high source
 shear.
\end{enumerate}
The first degree pair not certified automorphic by this sieve is
\((m,n)=(9,6)\).  At that pair the reduced data are
\[
 (a,b,d,t)=(3,2,3,3).
\]
\end{proposition}

\begin{proof}
Let \(f_m=[y^m]F\) and \(g_n=[y^n]G\).  The coefficient of
\(y^{m+n-1}\) in the Jacobian is
\[
 n f_m'g_n-m f_mg_n'=0.
\]
Write
\[
 d=\gcd(m,n),\qquad m=ad,\qquad n=bd,\qquad \gcd(a,b)=1.
\]
Unique factorization in \(\C[x]\), followed by constant rescaling, gives
\[
 f_m=h^a,\qquad g_n=h^b
\]
for a nonzero \(h\in\C[x]\).  Put
\[
 s=\deg h,\qquad t=\gcd(d,s),
\]
with \(\gcd(d,0)=d\).  The two partial-degree invariants are
\[
 \gcd(m,\deg f_m)=at,\qquad
 \gcd(n,\deg g_n)=bt.
\]

Choose \(L\) larger than the \(x\)-degree of every coefficient of
\(F,G\), and make the source shear \(y\mapsto y+x^L\).  The resulting
total degrees are
\[
 as+adL=at\,p,\qquad bs+bdL=bt\,p,
 \qquad p=\frac{s}{t}+\frac{d}{t}L.
\]
If \(s>0\), the two coefficients in the progression defining \(p\) are
coprime, so Dirichlet's theorem permits an arbitrarily large \(L\) for
which \(p\) is prime.  If \(s=0\), then \(t=d\), \(p=L\), and one
chooses \(L\) prime.  Consequently the total-degree gcd is \(tp\).

The cited total-degree results certify automorphy when \(t=1\) or
\(t=2\).  They also certify automorphy if either \(at\) or \(bt\) lies
in
\[
 \{1,4\}\cup\{\text{prime numbers}\},
\]
because the corresponding total degree is respectively a prime,
\(4p\), or a product of two primes.  These conclusions use one common
prime progression only.

For a minimal unresolved pair, a constant linear target change removes
\(m=n\).  If \(b=1\), subtracting a constant multiple of \(G^a\) from
\(F\) lowers the ordered \(y\)-degree pair.  Hence
\[
 a>b\geq2,\qquad \gcd(a,b)=1,\qquad t\geq3,\qquad t\mid d.
\]
The least possible tuple is
\[
 (a,b,t,d)=(3,2,3,3),
\]
which gives \((m,n)=(9,6)\).  Conversely, its partial invariants are
\((9,6)\), and the high-shear total-degree gcd is \(3p\); none belongs
to the cited classes.  Thus it is the first survivor of this sieve.
\end{proof}

\begin{remark}
Proposition~\ref{prop:orientation} does not assert that a Keller pair of
degree \((9,6)\) exists.  It is not used to prove
Theorem~\ref{thm:main}; it only motivates why that coefficient system
is examined.
\end{remark}

\section{Connected normalization and approximate roots}
\label{sec:normalization}

Assume for contradiction that \(h\) is not a cube in \(\C(x)\).  Put
\[
 L=\C(x)(H),\qquad H^3=h.
\]
Because \(\C\) contains the cube roots of unity, this is a connected
cyclic cubic extension.  Let \(\zeta\) be a primitive cube root of
unity and let its deck generator satisfy
\[
 \sigma(H)=\zeta H.
\]
With \(z=Hy\), define
\[
 f(x,z)=F(x,z/H),\qquad g(x,z)=G(x,z/H).
\]
Then \(f,g\in L[z]\) are monic of degrees nine and six and
\begin{equation}
 [f,g]_{x,z}=\frac{\lambda}{H}.
\label{eq:normalized-jac}
\end{equation}
The deck action sends \(z\) to \(\zeta z\) and fixes \(f,g\).

Translate \(z\) over \(L\) to kill the degree-five coefficient of
\(g\).  Its deck character makes the translated variable \(w\) satisfy
\(\sigma(w)=\zeta w\).  Write
\begin{equation}
 g=w^6+aw^4+bw^3+cw^2+dw+q.
\label{eq:g}
\end{equation}
The coefficient characters are:
\[
\begin{array}{c|ccccc}
 &a&b&c&d&q\\ \hline
 \wt&2&3&4&5&6\\
 \sigma&\zeta^2&1&\zeta&\zeta^2&1.
\end{array}
\]

\begin{lemma}[Triangular approximate-root form]
\label{lem:approx}
After a constant target shear and translation, there is a
\(j\in\C\) such that
\[
 \boxed{f=(g^{3/2})_+ +j(g^{1/2})_+.}
\]
Here \((\cdot)_+\) denotes the polynomial part of the distinguished
Laurent expansion at \(w=\infty\).
\end{lemma}

\begin{proof}
For \(0\leq k\leq9\), put
\[
 A_k=(g^{k/6})_+.
\]
Each \(A_k\) is monic of degree \(k\), so there is a unique triangular
expansion
\[
 f=A_9+\sum_{k=0}^8c_k(x)A_k,\qquad c_k\in L.
\]
The full Laurent series \(g^{k/6}\) commutes with \(g\), and its
negative part begins in degree \(-1\).  Therefore
\[
 \deg_w[A_k,g]\leq4.
\]
On the other hand,
\[
 [c_kA_k,g]=c_k'A_kg_w+c_k[A_k,g],
\]
and the first summand has leading term \(6c_k'w^{k+5}\).
Descending from \(k=8\) to \(k=0\), equation
\eqref{eq:normalized-jac} forces every \(c_k'=0\).  The constant field
of the algebraic function field \(L/\C\) is \(\C\), so \(c_k\in\C\).

The distinguished sixth root of \(g\) has leading term \(w\), and hence
\(\sigma(A_k)=\zeta^kA_k\).  Since \(f\) is invariant and the
coefficients \(c_k\) are constants, uniqueness of the triangular
expansion forces \(c_k=0\) unless \(3\mid k\).  A target shear removes
the \(c_6g\) term and a target translation removes \(c_0\).  Writing
\(c_3=j\) gives the result.
\end{proof}

Expanding Lemma~\ref{lem:approx} gives
\begin{align}
 f={}&w^9+\frac32aw^7+\frac32bw^6
 +\frac{3(a^2+4c)}8w^5+\frac{3(ab+2d)}4w^4 \notag\\
 &+\frac{-a^3+12ac+6b^2+16j+24q}{16}w^3
 -\frac{3(a^2b-4ad-4bc)}{16}w^2+L_1w+P,
\label{eq:f-expanded}
\end{align}
where
\begin{align}
 L_1={}&\frac{
 3a^4-24a^2c-24ab^2+64aj+96aq+96bd+48c^2
 }{128},\label{eq:L}\\
 P={}&\frac{
 3a^3b-6a^2d-12abc-2b^3+16bj+24bq+24cd
 }{32}.
\label{eq:P}
\end{align}
Assigning
\[
 \wt(a,b,c,d,q;j)=(2,3,4,5,6;6)
\]
makes every coefficient in \eqref{eq:f-expanded} weighted homogeneous
of the weight complementary to its power of \(w\).

\section{The four lower first integrals}
\label{sec:integrals}

Define
\begin{align}
I_{10}={3\over128}\bigl(&3a^5-24a^3c-36a^2b^2+32a^2j+48a^2q
 +96abd \notag\\
 &+48ac^2+48b^2c-128cj-192cq-96d^2\bigr),
\label{eq:I10}\\
I_{11}={3\over128}\bigl(&15a^4b-24a^3d-72a^2bc-24ab^3
 +64abj+96abq \notag\\
 &+96acd+48b^2d+48bc^2-128dj-192dq\bigr),
\label{eq:I11}\\
I_{12}=-{1\over512}\bigl(&15a^6-132a^4c-288a^3b^2
 +128a^3j+192a^3q \notag\\
 &+672a^2bd+336a^2c^2+768ab^2c-512acj-768acq
 -384ad^2 \notag\\
 &+72b^4-384b^2j-576b^2q-1152bcd-192c^3
 +1536jq+1152q^2\bigr),
\label{eq:I12}\\
I_{13}=-{1\over128}\bigl(&15a^5b-21a^4d-96a^3bc-48a^2b^3
 +64a^2bj+96a^2bq \notag\\
 &+120a^2cd+120ab^2d+144abc^2-64adj-96adq
 +48b^3c \notag\\
 &-128bcj-192bcq-96bd^2-144c^2d\bigr).
\label{eq:I13}
\end{align}

\begin{lemma}[Exact integration of the lower equations]
\label{lem:integrals}
Let
\[
 R_k=[w^k][f,g]_{x,w}
\]
after substituting \eqref{eq:g} and \eqref{eq:f-expanded}.  Then
\[
\boxed{
\begin{aligned}
 R_4&=I_{10}',\\
 R_3&=I_{11}',\\
 R_2&=I_{12}'+\frac12aR_4,\\
 R_1&=I_{13}'+\frac13bR_4+\frac13aR_3.
\end{aligned}}
\]
The terminal coefficient is
\[
 R_0=dP'-L_1q'.
\]
\end{lemma}

\begin{proof}
Differentiate \eqref{eq:g} and \eqref{eq:f-expanded}, form
\(f_xg_w-f_wg_x\), and collect coefficients of \(w\).
Substitution of \eqref{eq:I10}--\eqref{eq:I13} gives the displayed
four identities coefficient by coefficient.  The constant coefficient
is \(dP'-L_1q'\).  This is a polynomial identity over
\(\mathbb Q[a,b,c,d,q,j,a',b',c',d',q']\); no division or genericity
assumption is used.
\end{proof}

Equation \eqref{eq:normalized-jac} and
Lemma~\ref{lem:integrals} make all four \(I_r\) constant.  Their weights
give the deck characters
\[
 \sigma(I_r)=\zeta^rI_r.
\]
Connectedness and the constant field \(\C\) therefore imply
\begin{equation}
 \boxed{
 I_{10}=0,\qquad I_{11}=0,\qquad I_{12}=K,\qquad I_{13}=0
 }
\label{eq:levels}
\end{equation}
for one \(K\in\C\).  The terminal equation is
\begin{equation}
 \boxed{dP'-L_1q'=\frac{\lambda}{H}.}
\label{eq:terminal}
\end{equation}

\section{Exact decomposition of the level set}
\label{sec:decomposition}

Make the centered change
\[
 C=c-\frac{a^2}{4},\qquad
 D=d-\frac{ab}{2},\qquad
 Q=q-\frac{b^2}{4}.
\]
Direct substitution gives
\begin{align}
 I_{10}&={3\over8}
 (3C^2a-12CQ-8Cj-6D^2),\label{eq:center10}\\
 I_{11}&={3\over8}
 (3C^2b+6CDa-12DQ-8Dj),\label{eq:center11}\\
 I_{12}&={1\over8}\bigl(
 3C^3-3C^2a^2+18CDb+12CQa+8Caj \notag\\
 &\hspace{35mm}{}+6D^2a-18Q^2-24Qj\bigr),\label{eq:center12}\\
 I_{13}&={1\over16}\bigl(
 18C^2D-9C^2ab-6CDa^2+24CQb+16Cbj \notag\\
 &\hspace{35mm}{}+12D^2b+12DQa+8Daj\bigr).\label{eq:center13}
\end{align}
The key factorization is the exact identity
\begin{equation}
 \boxed{
 I_{13}=\frac98C^2D-\frac b3I_{10}-\frac a6I_{11}.
 }
\label{eq:keyfactor}
\end{equation}
On the levels \eqref{eq:levels}, it gives \(C^2D=0\).

\subsection{The branch \(C=0\)}

If \(C=0\), equation \eqref{eq:center10} gives \(D=0\), and
\[
 I_{12}=-\frac34Q(3Q+4j)=K.
\]
Thus \(Q\), being algebraic over the constant field, is constant.
Set
\[
 \rho=w^3+\frac a2w+\frac b2 .
\]
Substitution in \eqref{eq:g} and \eqref{eq:f-expanded} yields
\[
 g=\rho^2+Q,\qquad
 f=\rho^3+\left(\frac32Q+j\right)\rho.
\]
Both coordinates are functions of \(\rho\), so their Jacobian is zero,
contradicting \eqref{eq:normalized-jac}.  Hence \(C\neq0\).

\subsection{The branch \(C\neq0\)}

Now \eqref{eq:keyfactor} gives \(D=0\).
Equation \eqref{eq:center11} then gives \(b=0\), and
\eqref{eq:center10} gives
\begin{equation}
 a=\frac{4(3Q+2j)}{3C}.
\label{eq:a-centered}
\end{equation}
The remaining invariant is the cubic level
\begin{equation}
 \boxed{C^3=R(Q):=6Q^2+8jQ+\frac83K.}
\label{eq:cubic-level}
\end{equation}
Since \(b=D=0\), the original coefficients satisfy
\[
 d=0,\qquad q=Q,\qquad P=0.
\]
Reducing \eqref{eq:L} modulo \eqref{eq:cubic-level} gives
\begin{equation}
 L_1=\frac{S(Q)}{24C},\qquad
 S(Q):=21R(Q)+32(j^2-K).
\label{eq:S}
\end{equation}
The terminal equation becomes
\begin{equation}
 -L_1Q'=\frac{\lambda}{H}.
\label{eq:terminal-Q}
\end{equation}

\section{Exclusion of the two levels}
\label{sec:exclusion}

\subsection{The generic level \(K\neq j^2\)}

Both \(b\) and \(q\) are deck-invariant, so
\[
 Q=q-\frac{b^2}{4}\in\C(x).
\]
Cubing \eqref{eq:terminal-Q} and using \(C^3=R(Q)\) gives the
base-field identity
\begin{equation}
 \boxed{
 h=-24^3\lambda^3
 \frac{R(Q)}{S(Q)^3Q'^3}.
 }
\label{eq:h-rational}
\end{equation}
The discriminant of the quadratic \(S\) is
\[
 12096(j^2-K).
\]
Thus \(S\) has two distinct roots when \(K\neq j^2\).  Neither is a
root of \(R\), since a common root in \eqref{eq:S} would give
\(32(j^2-K)=0\).

The terminal equation makes \(Q\) nonconstant.  Regard it as a
rational map
\(\mathbb P^1_x\to\mathbb P^1_Q\).  At least one of the two distinct
finite roots \(\alpha\) of \(S\) has a finite preimage: the single point
\(x=\infty\) cannot lie above two different target values.  At such a
preimage \(x_0\), write
\[
 \ord_{x_0}(Q-\alpha)=e>0.
\]
Then
\[
 \ord_{x_0}S(Q)=e,\qquad
 \ord_{x_0}Q'=e-1,\qquad
 \ord_{x_0}R(Q)=0.
\]
Formula \eqref{eq:h-rational} gives \(h\) a pole of order
\(6e-3>0\) at the finite point \(x_0\), impossible for
\(h\in\C[x]\).  This excludes every generic level.

\subsection{The cuspidal level \(K=j^2\)}

Now \eqref{eq:cubic-level} is
\[
 C^3=6\left(Q+\frac{2j}{3}\right)^2.
\]
On \(C\neq0\), put
\[
 v=\frac{Q+2j/3}{C}.
\]
Then
\begin{equation}
 C=6v^2,\quad Q=6v^3-\frac{2j}{3},\quad
 a=4v,\quad c=10v^2,\quad b=d=0,
\label{eq:cusp-params}
\end{equation}
and the normalized pair becomes
\begin{align}
 f={}&w^9+6vw^7+21v^2w^5+35v^3w^3+\frac{63}{2}v^4w,
\label{eq:cusp-f}\\
 g={}&w^6+4vw^4+10v^2w^2+6v^3-\frac{2j}{3}.
\label{eq:cusp-g}
\end{align}
A direct calculation gives
\begin{equation}
 [f,g]_{v,w}=-567v^6.
\label{eq:cusp-jac}
\end{equation}

The depression used in Section~\ref{sec:normalization} occurred over
the function field, so polynomial descent of \(v\) cannot be assumed.
We recover exactly the polynomial quantities needed for the valuation
argument.

\begin{lemma}[Finite-map descent at the cusp]
\label{lem:finite}
After substitution back into the original variables,
\[
 v^3,\ vw\in\C[x,y].
\]
\end{lemma}

\begin{proof}
Put \(g_0=g+2j/3\).  The pair \((f,g_0)\) is positive
weighted-homogeneous for
\(\wt(v,w)=(2,1)\).  Its only common zero is the origin.
Indeed, on the coordinate axes this follows from
\[
 (f,g_0)|_{v=0}=(w^9,w^6),\qquad
 g_0|_{w=0}=6v^3.
\]
Away from the axes put \(T=w^2/v\).  A common nonzero zero would make
\[
\begin{aligned}
 T^3+4T^2+10T+6&=0,\\
 T^4+6T^3+21T^2+35T+\frac{63}{2}&=0,
\end{aligned}
\]
but the resultant of these two polynomials is \(567/8\).

Therefore \(f,g_0\) form a homogeneous system of parameters in the
positively graded ring \(\C[v,w]\).  It follows that
\(\C[v,w]\) is module-finite over \(\C[f,g_0]\), so \(v^3\) and \(vw\)
are integral over \(\C[f,g_0]\).

After returning to \(x,y\), the elements \(f,g_0\) are constant target
transforms of the original polynomial coordinates.  Hence
\(\C[f,g_0]\subset\C[x,y]\), and the same monic equations make
\(v^3,vw\) integral over \(\C[x,y]\).  Both are deck-invariant and
belong to the fixed field
\[
 L(y)^{\mu_3}=\C(x,y).
\]
Since \(\C[x,y]\) is normal, integrality now gives
\(v^3,vw\in\C[x,y]\).
\end{proof}

Write the depressed translation and cusp parameter as
\[
 w=H(y+s),\qquad v=\frac{V}{H},
 \qquad s,V\in\C(x).
\]
The character assertions justify this form: \(w/H\) and \(vH\) are
deck-invariant.  By Lemma~\ref{lem:finite},
\[
 vw=V(y+s)\in\C[x,y],
\]
so
\[
 V\in\C[x],\qquad Vs\in\C[x].
\]
Also
\[
 v^3=\frac{V^3}{h}\in\C[x,y]\cap\C(x)=\C[x],
\]
and therefore
\begin{equation}
 h\mid V^3.
\label{eq:divisibility}
\end{equation}

Combining \eqref{eq:normalized-jac},
\eqref{eq:cusp-jac}, and \(v=V/H\) gives
\begin{equation}
 \boxed{
 -189V^6(3hV'-Vh')=\lambda h^3.
 }
\label{eq:valuation-identity}
\end{equation}
Let \(x_0\) be a root of \(h\), and put
\[
 r=\ord_{x_0}V,\qquad s_0=\ord_{x_0}h.
\]
Divisibility \eqref{eq:divisibility} gives \(s_0\leq3r\).
If \(s_0\neq3r\), then the orders of the two sides of
\eqref{eq:valuation-identity} are
\[
 7r+s_0-1,\qquad 3s_0.
\]
Their equality gives \(7r-1=2s_0\).  Together with
\(s_0\leq3r\), this forces \(r=1,s_0=3\), contrary to
\(s_0\neq3r\).  Hence \(s_0=3r\) at every root of \(h\).

A root of \(V\) outside the zero set of \(h\) would give orders
\(7r-1>0\) and \(0\) in \eqref{eq:valuation-identity}, also
impossible.  Thus \(h\) and \(V\) have the same support, with
\[
 h=\gamma V^3,\qquad \gamma\in\C^\times.
\]
But then \(3hV'-Vh'=0\), contradicting
\eqref{eq:valuation-identity}.  This excludes the cuspidal level and
completes the proof of Theorem~\ref{thm:main}.

\section{The cube-\(h\) boundary}
\label{sec:boundary}

\begin{remark}[Why the proof stops]
If \(h\) is a cube in \(\C(x)\), factorization over \(\C\) gives
\(h=r^3\), up to a constant cube, for some \(r\in\C[x]\).  This
includes all nonzero constant \(h\).  Put \(z=ry\) and depress \(g\)
over \(\C(x)\).  The triangular part of Lemma~\ref{lem:approx} remains
valid, but there is no nontrivial deck character.  After removing the
constant \(g\)-multiple and target translation, its exact conclusion is
\[
\begin{aligned}
 f={}&(g^{3/2})_+
 +\kappa _8(g^{4/3})_+
 +\kappa _7(g^{7/6})_+
 +\kappa _5(g^{5/6})_+\\
 &+\kappa _4(g^{2/3})_+
 +j(g^{1/2})_+
 +\kappa _2(g^{1/3})_+
 +\kappa _1(g^{1/6})_+ ,
\end{aligned}
\]
with seven unrestricted constants at the upper-equation level.  In
particular, the \(w^8\)-coefficient need not vanish.  For instance,
\[
 F=(xy)^9+(xy)^8,\qquad G=(xy)^6
\]
when \(h=x^3\), and
\[
 F=y^9+y^8,\qquad G=y^6
\]
when \(h=1\), descend polynomially and retain that coefficient.  These
examples have zero Jacobian and are not Keller pairs; they only show
that descent and the upper equations do not remove the extra constants.
The lower cube-\(h\) system is outside Theorem~\ref{thm:main}.
\end{remark}

No Laurent-time reduction is used in this manuscript.  Such a reduction
requires additional hypotheses and is best treated separately rather
than being folded into the connected theorem.

\appendix
\section{Exact-computation companion}

The proof is algebraic.  A companion script in the research archive,
\begin{quote}\scriptsize
\path{current_context/verify_normal_degree_arithmetic_frontier_and_96_connected_reduction.py},
\end{quote}
performs exact symbolic checks of:
\begin{itemize}
 \item the arithmetic enumeration through \(m=9\);
 \item the approximate-root expansion
 \eqref{eq:f-expanded}--\eqref{eq:P};
 \item the four derivative identities in
 Lemma~\ref{lem:integrals};
 \item the centered formulas and factorization
 \eqref{eq:center10}--\eqref{eq:keyfactor};
 \item the generic-level discriminant and the cuspidal parametrization;
 \item the resultant \(567/8\), the cusp Jacobian
 \eqref{eq:cusp-jac}, and the descent identity
 \eqref{eq:valuation-identity}.
\end{itemize}
The script is a regression check against expansion and sign errors; it
does not replace the proof.

\paragraph{Disclosure.}
AI systems assisted with symbolic-search orchestration, adversarial
review, exact-code checking, and draft preparation.  The named author
is responsible for the mathematical claims, the source, and
independent verification before dissemination.

\begin{thebibliography}{9}

\bibitem{Moskowicz}
V. Moskowicz,
\emph{A variation on Magnus' theorem and its generalizations},
preprint, arXiv:1810.08202v2, 2018,
\url{https://arxiv.org/abs/1810.08202v2}.

\bibitem{GGHVApprox}
J.~A. Guccione, J.~J. Guccione, R. Horruitiner, and C. Valqui,
\emph{The Jacobian Conjecture: Approximate roots and intersection
numbers},
preprint, arXiv:1708.09367v2, 2018,
\url{https://arxiv.org/abs/1708.09367v2}.

\bibitem{GGVShape}
J.~A. Guccione, J.~J. Guccione, and C. Valqui,
\emph{On the shape of possible counterexamples to the Jacobian
Conjecture},
J. Algebra \textbf{471} (2017), 13--74;
preprint arXiv:1401.1784,
\url{https://doi.org/10.1016/j.jalgebra.2016.08.039}.

\bibitem{GGHVAlgorithms}
J.~A. Guccione, J.~J. Guccione, R. Horruitiner, and C. Valqui,
\emph{Some algorithms related to the Jacobian Conjecture},
preprint, arXiv:1708.07936, 2017,
\url{https://arxiv.org/abs/1708.07936}.

\bibitem{Karas}
M. Kara\'s,
\emph{On weighted bidegree of polynomial automorphisms of
\(\C^2\)},
Bull. Polish Acad. Sci. Math. \textbf{70} (2022), 107--114;
preprint arXiv:1201.3463,
\url{https://doi.org/10.4064/ba220430-21-3}.

\end{thebibliography}

\end{document}
